\makeatletter
\def\input@path{{./}{../}{../../}{../../../}{../../../../}{preamble/}{../preamble/}{../../preamble/}{../../../preamble/}}
\makeatother

\input{preamble_exams.tex}
\input{graphs/highlights.tex}
% \GTenHighlightsfalse % Uncomment to disable instructional highlights.
\input{graphs/specific_highlights.tex}
\COneTonePracticePalette

\input{graphs/uniform_linear_probability_distributions.tex}
\input{graphs/probability_distribution_graphs.tex}
\input{graphs/normal_distribution_graphs.tex}

\title{Grade 11 -- Term 1 \\ C1 Practice Activity 1 with Solutions}
\author{}
\date{}

\begin{document}
	\maketitle

	\setcounter{tocdepth}{1}
	\setcounter{secnumdepth}{2}
	\phantomsection
	\label{toc}
	\tableofcontents

	\section*{Evaluated criterion}
	\begin{itemize}
		\item C1: Characterises continuous random variables by distinguishing them from discrete variables and interpreting their probability density function, cumulative distribution function, and support in economic contexts.
	\end{itemize}


	\section{Introduction}
	\label{sec:introduction}
\Needspace{8\baselineskip}
\subsection*{Mathematical highlight guide}
\phantomsection\label{tone:C1:legend}
\PracticeHighlightLegend{G11-T1-C1}


	\SectionProblemsTableOfContents{
		\SectionProblemLink{sec:introduction-bernoulli}{Bernoulli Distribution} \\
		\SectionProblemLink{sec:introduction-binomial}{Binomial Distribution} \\
		\SectionProblemLink{sec:introduction-geometric}{Geometric Distribution} \\
		\SectionProblemLink{sec:introduction-poisson}{Poisson Distribution} \\
		\SectionProblemLink{sec:introduction-uniform}{Uniform Distribution} \\
		\SectionProblemLink{sec:introduction-triangular}{Triangular Distribution} \\
		\SectionProblemLink{sec:introduction-linear}{Linear Distribution} \\
		\SectionProblemLink{sec:introduction-normal}{Normal Distribution} \\
		\SectionProblemLink{sec:introduction-summary}{Summary}
	}\vspace{5mm}

	A \textbf{\ConeToneVariable{random} \ConeToneVariable{variable}} assigns a numerical value to the outcome of a random process. A discrete \ConeToneVariable{random} \ConeToneVariable{variable} takes separate, countable values (often counts such as $0,1,2,\ldots$). A continuous \ConeToneVariable{random} \ConeToneVariable{variable} can take any real value in an interval, such as a waiting time or price. Its \textbf{\ConeToneVariable{support}} is the set of possible values.

	For a discrete variable, the probability mass function (PMF) $p_X(\ConeToneVariable{x})=\ConeToneProbability{P}(\ConeToneVariable{X}=\ConeToneVariable{x})$ gives probability at each possible value, and the masses sum to $1$. For a continuous variable, the \CtwoToneDensity{probability} \CtwoToneDensity{density} \CtwoToneDensity{function} (PDF) $\CtwoToneDensity{f_X}(\ConeToneVariable{x})$ is a height, not a point probability: $\ConeToneProbability{P}(\ConeToneVariable{X}=\ConeToneVariable{x})=\ConeToneProbability{0}$, while interval probability is area, $\ConeToneProbability{P}(\ConeToneVariable{r}\le \ConeToneVariable{X}\le \CsevenToneSampleSD{s})=\int_r^s \CtwoToneDensity{f_X}(\ConeToneVariable{x})\,\ConeToneVariable{d}\ConeToneVariable{x}$. A valid density is nonnegative and has total area $1$. In both cases the \ConeToneCDF{cumulative} \ConeToneCDF{distribution} \ConeToneCDF{function} (CDF)
	\[
		\ConeToneCDF{F_X}(\ConeToneVariable{x})=\ConeToneProbability{P}(\ConeToneVariable{X}\le \ConeToneVariable{x})
	\]
	records the probability accumulated up to $\ConeToneVariable{x}$; it is a sum of PMF values or an integral of the PDF.

	\subsection{Bernoulli Distribution}
	\label{sec:introduction-bernoulli}
	A Bernoulli variable represents \emph{one} trial with two outcomes, coded success ($1$) and failure ($0$). Clues are ``one selected account,'' ``yes/no,'' or ``paid/defaulted.'' Its \ConeToneVariable{support} is $\{0,1\}$ and its parameter is the \ConeToneSuccess{success} \ConeToneSuccess{probability} $\ConeToneProbability{0}\le \ConeToneSuccess{p}\le\ConeToneProbability{1}$:
	\[
		\ConeToneProbability{P}(\ConeToneVariable{X}=\ConeToneVariable{x})=p^x(\ConeToneProbability{1}-\ConeToneSuccess{p})^{1-x},\quad \ConeToneVariable{x}\in\{0,1\},
		\qquad
		\ConeToneCDF{F_X}(\ConeToneVariable{x})=\begin{cases}\ConeToneCDF{0},&\ConeToneVariable{x}<0,\\1-\ConeToneSuccess{p},&0\le \ConeToneVariable{x}<1,\\1,&\ConeToneVariable{x}\ge1.\end{cases}
}\]
	The PMF assigns the two outcome probabilities; the CDF gives the probability of a coded outcome no greater than $\ConeToneVariable{x}$. For example, $\ConeToneVariable{X}=\ConeToneVariable{1}$ may mean that one randomly selected invoice is overdue.

	\secInicio[sec:introduction-bernoulli][sec:introduction]

	\subsection{Binomial Distribution}
	\label{sec:introduction-binomial}
	A binomial variable counts successes in a \emph{fixed} number $\CsixToneCount{n}$ of independent Bernoulli trials having the same \ConeToneSuccess{success} \ConeToneSuccess{probability} $\ConeToneSuccess{p}$. Its recognition clues are a stated \CsixToneCount{sample} \CsixToneCount{size}, binary outcomes, independence, and constant $\ConeToneSuccess{p}$. Its \ConeToneVariable{support} is $\{0,1,\ldots,\CsixToneCount{n}\}$ and its parameters are $\CsixToneCount{n}$ and $\ConeToneSuccess{p}$:
	\[
		\ConeToneProbability{P}(\ConeToneVariable{X}=\ConeToneVariable{x})=\binom{\CsixToneCount{n}}{\ConeToneVariable{x}}p^x(\ConeToneProbability{1}-\ConeToneSuccess{p})^{n-x},\quad \ConeToneVariable{x}=\ConeToneVariable{0},\ldots,\CsixToneCount{n},
		\qquad \ConeToneCDF{F_X}(\ConeToneVariable{x})=\sum_{k=0}^{\lfloor x\rfloor}\binom{\CsixToneCount{n}}{\CsixToneCount{k}}p^k(\ConeToneCDF{1}-\ConeToneSuccess{p})^{n-k},
	\]
	where the CDF is $0$ below $0$ and $1$ at or above $\CsixToneCount{n}$. The PMF gives the probability of exactly $\ConeToneVariable{x}$ successes; the CDF gives that of at most $\ConeToneVariable{x}$. An example is the number of renewed subscriptions among $\CsixToneCount{n}$ independently contacted customers.

	\secInicio[sec:introduction-binomial][sec:introduction]

	\subsection{Geometric Distribution}
	\label{sec:introduction-geometric}
	Here $\ConeToneVariable{X}$ counts the number of independent Bernoulli trials \emph{required to obtain the first success}; thus this activity uses \ConeToneVariable{support} $\{1,2,3,\ldots\}$. Look for repeated independent attempts, constant \ConeToneSuccess{success} \ConeToneSuccess{probability} $\ConeToneSuccess{p}$, and stopping at the first success:
	\[
		\ConeToneProbability{P}(\ConeToneVariable{X}=\ConeToneVariable{x})=(\ConeToneProbability{1}-\ConeToneSuccess{p})^{x-1}\ConeToneSuccess{p},\quad \ConeToneVariable{x}=\ConeToneVariable{1},2,\ldots,
		\qquad \ConeToneCDF{F_X}(\ConeToneVariable{x})=\begin{cases}\ConeToneCDF{0},&\ConeToneVariable{x}<1,\\1-(1-\ConeToneSuccess{p})^{\lfloor x\rfloor},&\ConeToneVariable{x}\ge1.\end{cases}
	\]
	The PMF describes $\ConeToneVariable{x}-1$ failures followed by a success; the CDF describes success occurring by trial $\ConeToneVariable{x}$. For example, $\ConeToneVariable{X}$ can be the number of sales pitches until the first order.

	\secInicio[sec:introduction-geometric][sec:introduction]

	\subsection{Poisson Distribution}
	\label{sec:introduction-poisson}
	A Poisson variable counts events in a fixed time, distance, area, or other interval when events occur separately at an approximately constant average rate $\ConeTonePoissonMean{\lambda}>0$. Clues are an event count, a fixed interval, and a stable rate. Its \ConeToneVariable{support} is $\{0,1,2,\ldots\}$:
	\[
		\ConeToneProbability{P}(\ConeToneVariable{X}=\ConeToneVariable{x})=e^{-\lambda}\frac{\ConeTonePoissonMean{\lambda^x}}{\ConeToneVariable{x}!},\quad \ConeToneVariable{x}=\ConeToneVariable{0},1,2,\ldots,
		\qquad \ConeToneCDF{F_X}(\ConeToneVariable{x})=e^{-\lambda}\sum_{k=0}^{\lfloor x\rfloor}\frac{\ConeTonePoissonMean{\lambda^k}}{\CsixToneCount{k}!}.
	\]
	The PMF gives the chance of exactly $\ConeToneVariable{x}$ events and the CDF the chance of no more than $\ConeToneVariable{x}$. A call centre might use it for the number of incoming calls in a fixed five-minute \CThreeTonePeriod{period}.

	\secInicio[sec:introduction-poisson][sec:introduction]

	\subsection{Uniform Distribution}
	\label{sec:introduction-uniform}
	A uniform variable treats all values in a bounded interval $[\ConeToneVariable{a},\ConeToneVariable{b}]$ as having equal density. A stated minimum and maximum plus ``equally plausible anywhere'' are decisive clues. Its \ConeToneVariable{support} is $[\ConeToneVariable{a},\ConeToneVariable{b}]$, with parameters $\ConeToneVariable{a}<\ConeToneVariable{b}$:
	\[
		\CtwoToneDensity{f_X}(\ConeToneVariable{x})=\begin{cases}\dfrac1{\ConeToneVariable{b}-\ConeToneVariable{a}},&\ConeToneVariable{a}\le \ConeToneVariable{x}\le \ConeToneVariable{b},\\\CtwoToneDensity{0},&\text{otherwise},\end{cases}
		\qquad
		\ConeToneCDF{F_X}(\ConeToneVariable{x})=\begin{cases}\ConeToneCDF{0},&\ConeToneVariable{x}<\ConeToneVariable{a},\\\dfrac{\ConeToneVariable{x}-\ConeToneVariable{a}}{\ConeToneVariable{b}-\ConeToneVariable{a}},&\ConeToneVariable{a}\le \ConeToneVariable{x}\le \ConeToneVariable{b},\\1,&\ConeToneVariable{x}>\ConeToneVariable{b}.\end{cases}
	\]
	The constant PDF makes equal-length subintervals equally probable; the CDF is the fraction of the \ConeToneVariable{support} lying at or below $\ConeToneVariable{x}$. A courier arrival time known only to be equally plausible from 2:00 to 2:30 p.m. is an example.

	\secInicio[sec:introduction-uniform][sec:introduction]

	\subsection{Triangular Distribution}
	\label{sec:introduction-triangular}
	A triangular variable has minimum $\ConeToneVariable{a}$, maximum $\ConeToneVariable{b}$, and \CtwoToneMode{mode} (most likely value) $\CtwoToneMode{c}$, where $\ConeToneVariable{a}<\CtwoToneMode{c}<\ConeToneVariable{b}$. Density rises linearly toward $\CtwoToneMode{c}$ and falls linearly afterward---the key distinction from one straight linear density across the whole \ConeToneVariable{support}. Its parameters are $\ConeToneVariable{a},\CtwoToneMode{c},\ConeToneVariable{b}$ and \ConeToneVariable{support} is $[\ConeToneVariable{a},\ConeToneVariable{b}]$:
	\[
	 \CtwoToneDensity{f_X}(\ConeToneVariable{x})=\begin{cases}
	 \CtwoToneDensity{0},&\ConeToneVariable{x}<\ConeToneVariable{a},\\
	 \dfrac{\ConeToneVariable{2}(\ConeToneVariable{x}-\ConeToneVariable{a})}{(\ConeToneVariable{b}-\ConeToneVariable{a})(\CtwoToneMode{c}-\ConeToneVariable{a})},&\ConeToneVariable{a}\le \ConeToneVariable{x}\le \CtwoToneMode{c},\\
	 \dfrac{\ConeToneVariable{2}(\ConeToneVariable{b}-\ConeToneVariable{x})}{(\ConeToneVariable{b}-\ConeToneVariable{a})(\ConeToneVariable{b}-\CtwoToneMode{c})},&\CtwoToneMode{c}<\ConeToneVariable{x}\le \ConeToneVariable{b},\\
	 0,&\ConeToneVariable{x}>\ConeToneVariable{b},
	 \end{cases}
	\]
	\[
	 \ConeToneCDF{F_X}(\ConeToneVariable{x})=\begin{cases}
	 \ConeToneCDF{0},&\ConeToneVariable{x}<\ConeToneVariable{a},\\
	 \dfrac{(\ConeToneVariable{x}-\ConeToneVariable{a})^2}{(\ConeToneVariable{b}-\ConeToneVariable{a})(\CtwoToneMode{c}-\ConeToneVariable{a})},&\ConeToneVariable{a}\le \ConeToneVariable{x}\le \CtwoToneMode{c},\\
	 1-\dfrac{(\ConeToneVariable{b}-\ConeToneVariable{x})^2}{(\ConeToneVariable{b}-\ConeToneVariable{a})(\ConeToneVariable{b}-\CtwoToneMode{c})},&\CtwoToneMode{c}<\ConeToneVariable{x}\le \ConeToneVariable{b},\\
	 1,&\ConeToneVariable{x}>\ConeToneVariable{b}.
	 \end{cases}
	\]
	The PDF expresses increasing then decreasing plausibility; its accumulated area is the CDF. A project cost with an expert minimum, maximum, and most likely estimate is a common business example.

	\secInicio[sec:introduction-triangular][sec:introduction]

	\subsection{Linear Distribution}
	\label{sec:introduction-linear}
	A linear distribution has a PDF $\CtwoToneDensity{f_X}(\ConeToneVariable{x})=m\ConeToneVariable{x}+\ConeToneVariable{d}$ throughout a bounded \ConeToneVariable{support} $[\ConeToneVariable{a},\ConeToneVariable{b}]$, and $0$ outside it. The constants must satisfy $m\ConeToneVariable{x}+\ConeToneVariable{d}\ge0$ on $[\ConeToneVariable{a},\ConeToneVariable{b}]$ and
	\[
		\int_a^b(m\ConeToneVariable{x}+\ConeToneVariable{d})\,\ConeToneVariable{d}\ConeToneVariable{x}=\frac{m}{2}(b^2-a^2)+\ConeToneVariable{d}(\ConeToneVariable{b}-\ConeToneVariable{a})=\ConeToneVariable{1}.
	\]
	Its \ConeToneVariable{support} is $[\ConeToneVariable{a},\ConeToneVariable{b}]$ and its parameters may be written $\ConeToneVariable{a},\ConeToneVariable{b},m,\ConeToneVariable{d}$ subject to these restrictions. The CDF is
	\[
	 \ConeToneCDF{F_X}(\ConeToneVariable{x})=\begin{cases}\ConeToneCDF{0},&\ConeToneVariable{x}<\ConeToneVariable{a},\\\dfrac{m}{2}(x^2-a^2)+\ConeToneVariable{d}(\ConeToneVariable{x}-\ConeToneVariable{a}),&\ConeToneVariable{a}\le \ConeToneVariable{x}\le \ConeToneVariable{b},\\1,&\ConeToneVariable{x}>\ConeToneVariable{b}.\end{cases}
	\]
	Clues specify one density line increasing or decreasing across the entire interval, without an interior \CtwoToneMode{mode}. The PDF gives relative concentration and the CDF its accumulated area. For instance, longer checkout times from 0 to 8 minutes may have density $\CtwoToneDensity{f}(\ConeToneVariable{t})=\ConeToneVariable{t}/\CtwoToneDensity{32}$, increasing linearly over the whole interval.

	\secInicio[sec:introduction-linear][sec:introduction]

	\subsection{Normal Distribution}
	\label{sec:introduction-normal}
	A normal variable models approximately symmetric, bell-shaped observations clustered around a centre, with values farther away increasingly uncommon. Its \ConeToneVariable{support} is all real numbers and its parameters are mean $\CfiveTonePopulationMean{\mu}$ and standard deviation $\CfiveTonePopulationSD{\sigma}>\CfiveTonePopulationSD{0}$:
	\[
		\CtwoToneDensity{f_X}(\ConeToneVariable{x})=\frac{\ConeToneVariable{1}}{\CfiveTonePopulationSD{\sigma}\sqrt{\CfiveTonePopulationVariance{2}\pi}}\exp\!\left[-\frac{(\ConeToneVariable{x}-\CfiveTonePopulationMean{\mu})^2}{2\CfiveTonePopulationVariance{\sigma^2}}\right],
		\qquad \ConeToneCDF{F_X}(\ConeToneVariable{x})=\ConeToneCDF{\Phi}\!\left(\frac{\ConeToneVariable{x}-\CfiveTonePopulationMean{\mu}}{\CfiveTonePopulationSD{\sigma}}\right).
	\]
	Here $\ConeToneCDF{\Phi}$ is the standard normal CDF (it is normally read from technology or a table). The PDF describes the bell's relative concentration and the CDF gives area to the left. A large firm's many employee commute times may be approximately normal when they are symmetric around a mean.

	\secInicio[sec:introduction-normal][sec:introduction]

	\subsection{Summary}
	\label{sec:introduction-summary}

	\noindent
	\begin{tabularx}{\textwidth}{|C{2.5cm}|C{2.8cm}|>{\raggedright\arraybackslash}X|}
		\hline
		\textbf{Type} & \textbf{Distribution} & \textbf{Typical situation} \\
		\hline
		Discrete & Bernoulli & One situation with two possible outcomes. \\
		\hline
		Discrete & Binomial & Counting how many times an outcome occurs across a fixed number of comparable cases. \\
		\hline
		Discrete & Geometric & Counting how many attempts are required until the first desired outcome. \\
		\hline
		Discrete & Poisson & Counting how many separate events occur during a fixed \CThreeTonePeriod{period} or exposure. \\
		\hline
		Continuous & Uniform & A bounded quantity for which no part of the possible range is expected more often than another. \\
		\hline
		Continuous & Triangular & A bounded quantity with a practical minimum, maximum, and one most common value. \\
		\hline
		Continuous & Linear & A bounded quantity whose observations become steadily more or less common across its range. \\
		\hline
		Continuous & Normal & Values concentrated around a typical value, with increasingly uncommon observations farther from it on either side. \\
		\hline
	\end{tabularx}

	\secInicio[sec:introduction-summary][sec:introduction]

	
\Needspace{8\baselineskip}
\subsection*{Notation reference}
\phantomsection\label{tone:C1:reference}
\COneToneDetailedReference{C1}
\section{Discrete Probabilities}
	\label{sec:discrete-probabilities}

	\SectionProblemsTableOfContents{
		\SectionProblemLink{sec:bogota-venture-decisions}{Bogotá Venture Decisions} \\
		\SectionProblemLink{sec:creative-digital-markets}{Creative and Digital Markets} \\
		\SectionProblemLink{sec:sustainable-trade-operations}{Sustainable Trade Operations} \\
		\SectionProblemLink{sec:finance-tourism-mobility}{Finance, Tourism, and Mobility}
	}

	% ----------------------------------------------------------------
	\subsection{Bogotá Venture Decisions}
	\label{sec:bogota-venture-decisions}

	Consider the following four situations.

	\begin{enumerate}
		\item A Bogotá fintech is monitoring one transfer to a freelance designer who needs the funds for this month's studio rent. The transfer will either clear during the banking window or be rejected, and no other result is recorded. The variable \(X_1\) records which of those two outcomes occurs so the platform can decide whether to release the funds or begin a review.
		\item A meal-preparation venture in Chapinero is calling offices to secure its first recurring lunch contract. Staff use the same current offer with one office after another and stop calling as soon as an office books a tasting; recent calling rounds show comparable responses from offices approached this way. The variable \(X_2\) records how many calls are needed through that first booking, which determines whether this sales channel is affordable.
		\item A flower exporter near Bogotá must decide whether to release a shipment without costly re-sorting. Inspectors select exactly 40 stems packed on the same line and record for each one whether it meets export grade; recent batches show a consistent pass experience across stems. The variable \(X_3\) counts the stems that pass, giving the manager evidence for the release decision.
		\item A Bogotá concert promoter must staff its customer-assistance team after announcing a schedule change. Refund requests are logged separately during one hour, and records from comparable announcement periods show activity spread fairly steadily through the hour rather than arriving on a timetable. The variable \(X_4\) counts the requests logged in that hour so the promoter can balance response times against labour cost.
	\end{enumerate}

	\textbf{Questions/Tasks.}
	\begin{enumerate}
		\CriterionTask
		{subsec:bogota-venture-decisions-c1}{1}
		{Characterise each random variable and assign it to the appropriate probability distribution.}
	\end{enumerate}

	\subsubsection*{C1 - Characterise the random variables.}
	\phantomsection
	\label{subsec:bogota-venture-decisions-c1}

	All four variables are discrete.
	\begin{enumerate}
		\item $\ConeToneVariable{X_1}$ is \textbf{Bernoulli}: one transfer is one trial with exactly two outcomes, clearing or rejection.
		\item $\ConeToneVariable{X_2}$ is \textbf{geometric}: independent, equal-probability calls continue until the first appointment.
		\item $\ConeToneVariable{X_3}$ is \textbf{binomial}: it counts passes in a fixed set of $40$ independent binary inspections with constant pass probability.
		\item $\ConeToneVariable{X_4}$ is \textbf{Poisson}: it counts separate events in a fixed hour under an approximately stable rate.
	\end{enumerate}

	\secInicio[sec:bogota-venture-decisions][sec:discrete-probabilities]

	% ----------------------------------------------------------------
	\subsection{Creative and Digital Markets}
	\label{sec:creative-digital-markets}

	Consider the following four situations.

	\begin{enumerate}
		\item A Bogotá podcast studio runs a two-day sponsorship campaign for a series about Colombian creative businesses. Enquiries appear as separate messages throughout comparable campaigns, with no recurring surge during the window and no scheduled sender list. The variable \(X_1\) counts the enquiries received in those two days so the studio can staff replies and judge whether the campaign is commercially worthwhile.
		\item A social-commerce seller in Suba is testing 25 paid advertisements before committing more working capital. Each advertisement uses the same offer with a comparable audience, and the seller records whether it leads to a purchase during the campaign window; past tests under these conditions have produced similar response patterns. The variable \(X_2\) counts the advertisements that lead to purchases so she can decide whether to expand the campaign.
		\item A Colombian streaming startup needs its first paying subscriber from a new visitor segment before expanding its promotion. It presents the same price to successive eligible visitors and ends the test when someone first subscribes; recent visitors in the segment have responded comparably to this offer. The variable \(X_3\) records how many offers are presented through that first subscription, informing the cash needed to enter the segment.
		\item An independent Bogotá animator submits one quote for an advertising project from an overseas client. The quote will either be accepted, filling the production calendar with foreign-currency income, or declined, leaving the month open for local commissions. The variable \(X_4\) records which outcome occurs so the animator can schedule work and plan cash flow.
	\end{enumerate}

	\textbf{Questions/Tasks.}
	\begin{enumerate}
		\CriterionTask
		{subsec:creative-digital-markets-c1}{1}
		{Characterise each random variable and assign it to the appropriate probability distribution.}
	\end{enumerate}

	\subsubsection*{C1 - Characterise the random variables.}
	\phantomsection
	\label{subsec:creative-digital-markets-c1}

	All four variables are discrete.
	\begin{enumerate}
		\item $\ConeToneVariable{X_1}$ is \textbf{Poisson}: it counts separate enquiries within a fixed two-day interval at an approximately stable rate.
		\item $\ConeToneVariable{X_2}$ is \textbf{binomial}: it counts purchases in exactly $25$ independent binary advertising trials with constant \ConeToneSuccess{success} \ConeToneSuccess{probability}.
		\item $\ConeToneVariable{X_3}$ is \textbf{geometric}: it counts independent offers through the first subscription when the \ConeToneSuccess{success} \ConeToneSuccess{probability} is constant.
		\item $\ConeToneVariable{X_4}$ is \textbf{Bernoulli}: a single quote has exactly two possible outcomes, accepted or declined.
	\end{enumerate}

	\secInicio[sec:creative-digital-markets][sec:discrete-probabilities]

	% ----------------------------------------------------------------
	\subsection{Sustainable Trade Operations}
	\label{sec:sustainable-trade-operations}

	Consider the following four situations.

	\begin{enumerate}
		\item A Bogotá sustainable-fashion brand inspects one returned jacket before settling the customer's refund. The inspection will classify it either as repairable for resale or as unsuitable for repair and destined for recycling. The variable \(X_1\) records that classification because it changes both the refund reserve and the value recoverable from inventory.
		\item A Colombian cacao exporter reviews a batch of 100 customs declarations before dispatch. Errors are discovered as separate omissions or mismatches, and in recent batches prepared by the same process they have appeared sporadically across the paperwork without a predictable position. The variable \(X_2\) counts the errors in the batch so the exporter can contract enough specialist review and avoid costly delays.
		\item A Colombian crafts cooperative sends catalogues to exactly 18 international retailers to plan an export launch. All receive the same wholesale terms, and for each retailer the cooperative records whether a sample is requested; comparable retailers contacted in recent rounds have shown similar responses. The variable \(X_3\) counts the sample requests so the cooperative can buy materials and reserve artisan time without overproducing.
		\item A Bogotá solar-equipment importer reviews prospective suppliers until it finds the first one meeting its certification rule. The same due-diligence process is applied to one candidate after another, and recent candidates drawn from the approved trade directory have shown comparable qualification records. The variable \(X_4\) records how many suppliers are checked through the first qualifying one, helping forecast sourcing time and review costs.
	\end{enumerate}

	\textbf{Questions/Tasks.}
	\begin{enumerate}
		\CriterionTask
		{subsec:sustainable-trade-operations-c1}{1}
		{Characterise each random variable and assign it to the appropriate probability distribution.}
	\end{enumerate}

	\subsubsection*{C1 - Characterise the random variables.}
	\phantomsection
	\label{subsec:sustainable-trade-operations-c1}

	All four variables are discrete.
	\begin{enumerate}
		\item $\ConeToneVariable{X_1}$ is \textbf{Bernoulli}: one jacket produces exactly one of two outcomes, repairable or not repairable.
		\item $\ConeToneVariable{X_2}$ is \textbf{Poisson}: it counts separate errors in a fixed batch under an approximately stable rate.
		\item $\ConeToneVariable{X_3}$ is \textbf{binomial}: it counts sample requests across $18$ fixed, independent binary trials with equal \ConeToneSuccess{success} \ConeToneSuccess{probability}.
		\item $\ConeToneVariable{X_4}$ is \textbf{geometric}: independent supplier checks continue until the first qualifying supplier, with constant \ConeToneSuccess{success} \ConeToneSuccess{probability}.
	\end{enumerate}

	\secInicio[sec:sustainable-trade-operations][sec:discrete-probabilities]

	% ----------------------------------------------------------------
	\subsection{Finance, Tourism, and Mobility}
	\label{sec:finance-tourism-mobility}

	Consider the following four situations.

	\begin{enumerate}
		\item A digital lender reviews Colombian microbusiness applications until it approves the first loan for a pilot product. Analysts apply the same lending rule to successive applicants from one eligibility list and stop after the first approval; earlier lists assembled in this way have shown comparable approval experience. The variable \(X_1\) records how many applications are reviewed through that approval so the lender can budget underwriting time and acquisition cost.
		\item A Bogotá bicycle-share operator monitors breakdown reports on a fixed 20-kilometre network segment for one week. Each report concerns a separate incident, and operational records for ordinary weeks show incidents scattered through the segment and week without a recurring timetable. The variable \(X_2\) counts the reports in that period so the operator can position mechanics and spare parts without excess inventory.
		\item A Colombian tourism platform presents the same local-experience package to exactly 60 eligible travellers. For every traveller it records whether the package is purchased, and previous groups recruited through the same channel have responded comparably. The variable \(X_3\) counts the buyers so the platform can reserve enough Bogotá guides without making unnecessary advance payments.
		\item A currency-management app monitors one scheduled dollar payment for a Colombian importer buying production inputs. By the deadline, the payment will either have executed or remain unexecuted; those are the only operational outcomes used by the app. The variable \(X_4\) records which occurs because a missed deadline could interrupt the shipment and require costly short-term finance.
	\end{enumerate}

	\textbf{Questions/Tasks.}
	\begin{enumerate}
		\CriterionTask
		{subsec:finance-tourism-mobility-c1}{1}
		{Characterise each random variable and assign it to the appropriate probability distribution.}
	\end{enumerate}

	\subsubsection*{C1 - Characterise the random variables.}
	\phantomsection
	\label{subsec:finance-tourism-mobility-c1}

	All four variables are discrete.
	\begin{enumerate}
		\item $\ConeToneVariable{X_1}$ is \textbf{geometric}: equal-probability independent reviews continue until the first approval.
		\item $\ConeToneVariable{X_2}$ is \textbf{Poisson}: it counts separate reports over a fixed network segment and week at an approximately stable rate.
		\item $\ConeToneVariable{X_3}$ is \textbf{binomial}: it counts buyers in exactly $60$ independent binary trials with constant purchase probability.
		\item $\ConeToneVariable{X_4}$ is \textbf{Bernoulli}: one scheduled payment has exactly two outcomes, on time or not on time.
	\end{enumerate}

	\secInicio[sec:finance-tourism-mobility][sec:discrete-probabilities]


	\section{Continuous Probabilities}
	\label{sec:continuous-probabilities}

	\SectionProblemsTableOfContents{
		\SectionProblemLink{sec:events-exports-platforms}{Events, Exports, and Platforms} \\
		\SectionProblemLink{sec:food-finance-content}{Food, Finance, and Content} \\
		\SectionProblemLink{sec:tourism-logistics-design}{Tourism, Logistics, and Design} \\
		\SectionProblemLink{sec:technology-commerce-sustainability}{Technology, Commerce, and Sustainability}
	}

	% ----------------------------------------------------------------
	\subsection{Events, Exports, and Platforms}
	\label{sec:events-exports-platforms}

	Consider the following four situations.

	\begin{enumerate}
		\item A Bogotá concert organiser must fit an artist's sound check between venue setup and opening the doors. Production records place the duration \(X_1\) from 30 to 70 minutes, with 45 minutes occurring most often; durations become less frequent as they approach either operational extreme. This pattern helps the organiser reserve enough venue time without routinely paying unnecessary overtime.
		\item A flower exporter near Bogotá purchases sleeves and cartons before filling overseas orders. On its mature production line, stem length \(X_2\) clusters around 60 centimetres, with roughly balanced departures above and below that typical length and progressively fewer stems far from it. The exporter uses this behavior to select packaging sizes that protect saleable output and limit waste.
		\item A Bogotá ride platform is testing promotional credits while keeping each offer within COP~5,000 to COP~15,000. For the credit amount \(X_3\), past randomised campaigns give no reason to expect one part of that range more often than another. Sampling offers in this way lets the platform compare demand across the budget range before choosing a standard incentive.
		\item A marketplace connecting Colombian video creators with advertisers must budget one short-video placement priced from COP~100,000 to COP~500,000. Its transaction history shows that the observed price \(X_4\) becomes steadily more common as the amount rises across that range, without a reversal before the upper limit. The marketplace uses this pattern to set campaign deposits and avoid a working-capital shortfall.
	\end{enumerate}

	\textbf{Questions/Tasks.}
	\begin{enumerate}
		\CriterionTask
		{subsec:events-exports-platforms-c1}{1}
		{Characterise each random variable and assign it to the appropriate probability distribution.}
	\end{enumerate}

	\subsubsection*{C1 - Characterise the random variables.}
	\phantomsection
	\label{subsec:events-exports-platforms-c1}

	All four variables are continuous.
	\begin{enumerate}
		\item $\ConeToneVariable{X_1}$ is \textbf{triangular}: it has bounded minimum and maximum values, an interior most likely value, and density that rises then falls linearly.
		\item $\ConeToneVariable{X_2}$ is \textbf{normal}: lengths are approximately symmetric around a centre and become progressively rarer with distance from it.
		\item $\ConeToneVariable{X_3}$ is \textbf{uniform}: it is bounded and every value in the interval is equally plausible.
		\item $\ConeToneVariable{X_4}$ is \textbf{linear}: its density is one increasing straight line across the full bounded \ConeToneVariable{support}.
	\end{enumerate}

	\secInicio[sec:events-exports-platforms][sec:continuous-probabilities]

	% ----------------------------------------------------------------
	\subsection{Food, Finance, and Content}
	\label{sec:food-finance-content}

	Consider the following four situations.

	\begin{enumerate}
		\item An artisan bakery supplying Bogotá restaurants monitors loaf weight \(X_1\) because portions affect ingredient margins and clients' menu costs. Mature production records concentrate near 800 grams, with similar variation on either side and increasingly few loaves at large departures from that amount. The bakery uses this pattern to plan flour purchases and reduce waste without under-serving clients.
		\item An independent worker in Engativá uses a finance app to move part of each day's earnings into savings between 8:00 and 10:00 p.m. For the transfer time \(X_2\), processing history offers no reason to anticipate one part of the two-hour window more often than another. The worker relies on that window to keep funds available for evening expenses while the app spreads its workload.
		\item A Bogotá video creator is preparing a fixed-price proposal whose margin depends on editing effort and collaborator fees. The total editing cost \(X_3\) cannot practically fall below COP~900,000 or exceed COP~1,800,000, and similar projects most often cost COP~1,200,000; costs appear less often as they approach either limit. The creator uses these estimates to quote a fee that protects the project margin.
		\item A Colombian digital bank studies identity-check duration \(X_4\) to size servers without overspending on capacity. Verifications take from zero to five minutes in its operating records, and durations become steadily less common from the shortest end to the longest, with no change in that downward pattern. The bank uses the behavior to plan infrastructure while controlling account-opening delays.
	\end{enumerate}

	\textbf{Questions/Tasks.}
	\begin{enumerate}
		\CriterionTask
		{subsec:food-finance-content-c1}{1}
		{Characterise each random variable and assign it to the appropriate probability distribution.}
	\end{enumerate}

	\subsubsection*{C1 - Characterise the random variables.}
	\phantomsection
	\label{subsec:food-finance-content-c1}

	All four variables are continuous.
	\begin{enumerate}
		\item $\ConeToneVariable{X_1}$ is \textbf{normal}: weights are approximately symmetric about a centre and taper in frequency on both sides.
		\item $\ConeToneVariable{X_2}$ is \textbf{uniform}: every instant in a bounded two-hour interval is equally plausible.
		\item $\ConeToneVariable{X_3}$ is \textbf{triangular}: it has stated bounds and an interior \CtwoToneMode{mode}, with linear rise and fall in density.
		\item $\ConeToneVariable{X_4}$ is \textbf{linear}: the density decreases in one straight line throughout its bounded \ConeToneVariable{support}.
	\end{enumerate}

	\secInicio[sec:food-finance-content][sec:continuous-probabilities]

	% ----------------------------------------------------------------
	\subsection{Tourism, Logistics, and Design}
	\label{sec:tourism-logistics-design}

	Consider the following four situations.

	\begin{enumerate}
		\item A Colombian craft brand must finance materials and artisan wages before an overseas buyer pays for the next order. Recent small export orders range from COP~1 million to COP~11 million, and the observed order value \(X_1\) becomes steadily less common as value rises across that range, without reversing direction. The brand uses this behavior to arrange working capital while avoiding an unnecessarily large credit line.
		\item A La Candelaria tour operator coordinates guides with timed reservations at Bogotá food businesses. Route duration \(X_2\) has a practical minimum of two hours and maximum of four hours, and three hours occurs most often; observations thin out toward both limits. The operator uses this experience to protect reservations from delays without routinely paying guides for excessive standby time.
		\item A logistics broker handling an import through Colombia must dispatch a truck after customs release but before a reserved warehouse slot. Release time \(X_3\) will occur between 9:00 a.m. and noon, and comparable clearances give the broker no reason to favour any part of that window. This information guides a dispatch choice that limits both truck waiting charges and late-arrival penalties.
		\item A Bogotá footwear designer is choosing a viable width range before funding a Colombian production run. In a large target-market sample, foot width \(X_4\) is concentrated near a typical measurement, with comparable variation on either side and progressively fewer observations at more distant measurements. The designer uses that pattern to cover most buyers without carrying too many costly variants.
	\end{enumerate}

	\textbf{Questions/Tasks.}
	\begin{enumerate}
		\CriterionTask
		{subsec:tourism-logistics-design-c1}{1}
		{Characterise each random variable and assign it to the appropriate probability distribution.}
	\end{enumerate}

	\subsubsection*{C1 - Characterise the random variables.}
	\phantomsection
	\label{subsec:tourism-logistics-design-c1}

	All four variables are continuous.
	\begin{enumerate}
		\item $\ConeToneVariable{X_1}$ is \textbf{linear}: its bounded-\ConeToneVariable{support} density decreases along one straight line from one endpoint to the other.
		\item $\ConeToneVariable{X_2}$ is \textbf{triangular}: the bounded duration has a most likely interior value and two linear density segments.
		\item $\ConeToneVariable{X_3}$ is \textbf{uniform}: every time in the stated bounded window is equally plausible.
		\item $\ConeToneVariable{X_4}$ is \textbf{normal}: measurements are approximately symmetric around a centre and less frequent farther from it.
	\end{enumerate}

	\secInicio[sec:tourism-logistics-design][sec:continuous-probabilities]

	% ----------------------------------------------------------------
	\subsection{Technology, Commerce, and Sustainability}
	\label{sec:technology-commerce-sustainability}

	Consider the following four situations.

	\begin{enumerate}
		\item A Bogotá cloud startup must reserve infrastructure for a client whose storage use can range from zero to 20 terabytes. Contract records show that the observed demand \(X_1\) becomes steadily more common as use rises across that range, without turning downward before the contractual cap. The startup uses this behavior to buy enough capacity while limiting the cost of unused resources.
		\item A circular-economy venture in Kennedy bids for used appliances while preserving margin for parts, labour, and warranties. The resale price \(X_2\) has a practical minimum of COP~400,000 and maximum of COP~900,000, with COP~650,000 most common in the venture's records; prices occur less often nearer either extreme. This experience helps the buyer set a prudent inventory bid.
		\item An international e-commerce platform adjusts Colombian sellers' peso receipts for exchange-rate movements. For the adjustment \(X_3\), treasury policy limits the change to between negative 1.5 percent and positive 1.5 percent, and historical transactions provide no reason to expect one part of that range more often than another. Sellers use this information to test margins and choose an adequate currency-risk buffer.
		\item A Colombian coffee exporter monitors moisture percentage \(X_4\) because unsuitable storage can reduce inventory value before shipment. Readings from bags handled under the usual process gather around 11 percent, show similar behavior below and above that typical level, and become increasingly scarce farther away. The exporter uses the pattern to prioritise inspections and protect export quality.
	\end{enumerate}

	\textbf{Questions/Tasks.}
	\begin{enumerate}
		\CriterionTask
		{subsec:technology-commerce-sustainability-c1}{1}
		{Characterise each random variable and assign it to the appropriate probability distribution.}
	\end{enumerate}

	\subsubsection*{C1 - Characterise the random variables.}
	\phantomsection
	\label{subsec:technology-commerce-sustainability-c1}

	All four variables are continuous.
	\begin{enumerate}
		\item $\ConeToneVariable{X_1}$ is \textbf{linear}: its density increases along one straight line across a bounded \ConeToneVariable{support}.
		\item $\ConeToneVariable{X_2}$ is \textbf{triangular}: stated bounds, an interior most likely price, and linear rise and fall identify it.
		\item $\ConeToneVariable{X_3}$ is \textbf{uniform}: all real values in a bounded interval are equally plausible.
		\item $\ConeToneVariable{X_4}$ is \textbf{normal}: moisture values are approximately symmetric around a central value and rarer farther away.
	\end{enumerate}

	\secInicio[sec:technology-commerce-sustainability][sec:continuous-probabilities]


\end{document}
